The relation between flux $F$ and luminosity $L$:
$F = \dfrac{L}{4\pi d^2}$.

For Vega:
\[
F_V = \frac{L_V}{4\pi\,d_V^2}
\]
For Supernova:
\[
F_{SN} = \frac{L_{SN}}{4\pi\,d_{SN}^2}
\]
\begin{align*}
F_{SN} = 1.6\times10^{-7}F_V
\quad\Longrightarrow\quad
\frac{L_{SN}}{4\pi\,d_{SN}^2}
  &= \left(1.6\times10^{-7}\right)\frac{L_V}{4\pi\,d_V^2}\\
L_{SN}\,d_V^2 &= \left(1.6\times10^{-7}\right)L_V\,d_{SN}^2
\end{align*}
\hfill(25\%)
\begin{align*}
d_{SN} &= d_V\left[\frac{1}{1.6\times10^{-7}}\,
  \frac{L_{SN}}{L_V}\right]^{1/2}
= 7.76\left[\frac{1}{1.6\times10^{-7}}\,
  \frac{5.8\times10^{9}}{130}\right]^{1/2}
= 129\,581\,812.59\ \mathrm{pc}\\
&= 1.3\times10^{8}\ \mathrm{pc}
\end{align*}
\hfill(35\%)

If $t_H$ is the Hubble time, then the relation between redshift, distance, and
Hubble time is:
\[
t_H = \frac{d_{SN}}{cz}
= \left[\frac{1.3\times10^{8}\ \mathrm{pc}}
  {(1\ \mathrm{lyr/yr})(0.03)}\right]
  \left(\frac{3.26\ \mathrm{lyr}}{1\ \mathrm{pc}}\right)
= 1.41\times10^{10}\ \mathrm{years}
\]
\hfill(40\%)
